\section{Introducción específica}
\subsection{Descripción del equipo}
El dispositivo dispone de dos botones, uno de los cuales permite seleccionar el modo de funcionamiento.
Esta configuración es necesaria debido a que los valores normales de saturación de oxígeno en sangre, frecuencia respratoria y frecuencia cardíaca difieren según la edad del usuario.
Con esta consideración, se cuentan 2 modos de funcionamiento: adulto o pediátrico. 
Una vez establecidos los parametros correspondientes al modo seleccionado, el sensor óptico MAX30102 se coloca en el dedo del usuario.

\begin{figure}[h]
    \centering
    \includegraphics[height=0.3 \textwidth]{ppg}
    \hspace{10pt}
    \includegraphics[height=0.3 \textwidth]{ppg2}
    \caption{Funcionamiento del sensor MAX30102.}
    \label{fig:ppg2}
\end{figure}

Este sensor se utiliza para la medición de señales de fotopletismografía (PPG), el mismo integra dos dfuentes de luz LED (una roja  y una infrarroja), un fotodiodo y un sistema de conversión analógico-digita.
Debido a que el volumen de sangre en los vasos varía con cada latido cardíaco, la cantidad de luz absorbida también cambia de forma periódica.
Las variaciones de luz absorbida son captadas por el fotodiodo del sensor, el cual convierte la intensidad de la luz reflejada en una señal eléctrica.
Los valores digitalizados corresponden a las mediciones de luz detectada por los LEDs rojo e infrarrojo, los cuales constituyen la señal PPG.
Estos datos son transmitidos al microcontrolador mediante $I^2C$.
Posteriormente, las señales adquiridas son sometidas a una etapa de procesamiento digital, en la cual se aplican los algoritmos necesarios para calcular los valores de saturación de oxígeno en sangre, frecuencia cardíaca y frecuencia respiratoria.
% A continuación se detalla brevemente dicho procesamiento: 

Para la estimación de la frecuencia cardíaca, se utiliza principalmente la señal PPG infrarroja.
Se aplica un filtrado pasa-banda en el rango aproximado de 0.5 a 5Hz, con el objetivo de eliminar ruidos para posteriormente detectar los máximos locales de la señal filtrada, que corresponden a latidos válidos.
A partir del intervalo temporal entre picos consecutivos, se calcula la frecuencia cardíaca en latidos por minuto.
En el caso de la saturación de oxígeno en sangre ($SpO_{2}$), se utilizan ambas señales PPG. Cada una de estas señales se descompone en dos componentes, una componente continua (DC) y una componente alterna (AC). A partir de estas componentes se calcula una relación entre las señales roja e infrerroja.

\begin{equation}
    R = \frac{\frac{AC_{RED}}{DC_{RED}}}{\frac{AC_{IR}}{DC_{IR}}}
\end{equation}

\begin{equation}
    SpO_{2} \approx 110 -25*R
\end{equation}

Por último, la frecuencia respiratoria puede estimarse analizando las oscilaciones de baja frecuencia presentes en la señal infrarroja.
Se aplica un filtro pasa-bajos con una frecuencia cercana a 0.4 Hz.
Finalmente, mediante la detección de picos o el análisis de la frecuencia dominante de la señal filtrada, se obtiene el número de ciclos respiratorios por minuto. 
Luego de la etapa de procesamiento, y considerando los rangos previamente establecidos según el modo de funcionamiento seleccionado, el sistema transmite en tiempo real la información obtenida tanto al display como a un dispositivo movil en caso de haberse establecido la conexión mediante bluetooth.
Los datos transmitidos incluyen : 

\begin{itemize}
    \item Detección de apnea (presencia o ausencia). 
    \item Frecuencia respiratoria actual. 
    \item Frecuencia cardíaca actual. 
    \item Saturacióin de oxígeno en sangre actual.
\end{itemize}

En caso de que la frecuencia respiratoria descienda por debajo de los límites considerados peligrosos, el sistema activa el buzzer y un led para alertar al usuario.
Esta alarma puede desactivarse manualmente mediante el botón de apagado correpondiente.
Asimismo, el dispositivo incorpora un modo de bajo consumo energético, el cual se activa automáticamente cuando no se detecta contacto del usuario con el sensor durante un período prolongado de tiempo.

\subsection{Requerimientos funcionales}

\begin{table}[h]
    \centering
    \hspace*{-.8cm}
    % se podría quitar la columna ID
    \begin{tabular}{|c|c|c|c|}
        \hline
        Grupo & ID & Requerimiento Funcional & Descripción \\ 
        \hline
        Adquisición& 1.1 & Adquisición de señales & Capturar señales ópticas ROJO/IR del MAX30102 via I2C.\\ \hline
         & 1.2 &Lectura periódica& Obtener muestras desde la FIFO del sensor.\\ \hline 
         &1.3 & Análisis de variación& Detectar variaciones innecesarias para $SpO_{2}$ y respiración \\ \hline
        Procesamiento& 2.1 & Filtrado digital& Aplicar filtrado \\ \hline
        & 2.2 &Cálculo de $SpO_{2}$ & Calcular $SpO_{2}$ mediante el ratio R\\ \hline
        & 2.3 & Cálculo de respiración & Estimar frecuencia respiratoria mediante\\ \hline
        & 2.3 & Cálculo de respiración & análisis de envolvente o línea base. \\ \hline
         & 2.4 & Empaquetado de datos & Generar paquete con $SpO_{2}$ , respiración y estado. \\ \hline
        Indicadores& 3.1 & Alarma crítica & Activar alarma sonora cuando se detecta apnea. \\ \hline
         & 3.2 & Indicadores LED & Mostrar estados (crítico, normal, error) mediante LEDs. \\ \hline
        Comunicación & 4.1 & Bluetooth & Transmitir datos vía HM-10 por UART. \\ \hline
        & 4.2 & Envío de parámetros & Enviar $SpO_{2}$, respiraciones y estado del paciente a la app. \\ \hline
         & 4.3 & Registro de eventos & Informar alarmas, desconexiones o estados especiales. \\ \hline
        & 4.4 & Display & Mostrar parámetros de interes por pantalla.\\ \hline
         % & 4.5 & Recepción de & Recibir umbrales y parámetros desde la aplicación. \\
         % &  & configuración & Recibir umbrales y parámetros desde la aplicación. \\ \hline
        Alimentación& 5.1 & Alimentación & Alimentación vía USB \\ \hline
         & 5.2 & Bajo consumo & Entrar en modo Sleep cuando no hay actividad. \\ \hline
        & 5.3 & Reactivación & Retomar actividad ante movimiento o comando. \\ \hline
    \end{tabular}
    \caption{Tabla de requerimientos}
    \label{tab:tab_requerimientos}
\end{table}

\subsection{Casos de uso}
\subsection{Descripción de módulos externos utilizados}

\begin{itemize}
    \item Display LCD 16x2
    \item Sensor MAX30102
    \item Módulo Bluetooth HM-10
\end{itemize}

\subsubsection{Display LCD 16x2}
El display utilizado es un display LCD de 16x2 como se puede observar en la figura \ref{fig:LCD}.
Dichas conexiones se realizaron con una configuración de 4 bits.
El mismo efectúa la comunicación con la placa mediante pines \textit{GPIO}. 

\begin{figure} [H]
    \centering
    \includegraphics[width=0.5\linewidth]{lcd_display.png}
    \caption{Módulo 1602A}
    \label{fig:LCD}
\end{figure}

\subsubsection{HM-10}
Para la transmisión de los datos hacia el dispositivo móvil del usuario se emplea un módulo Bluetooth HM-10.
La comunicación entre el módulo Bluetooth y la placa de desarrollo se realiza usando la interfaz UART.
A través de esta interfaz serie, el microcontrolador transmite al módulo los parámetros de interés para su posterior envío al dispositivo móvil del usuario.

\begin{figure}[H]
    \centering
    \includegraphics[width=0.5\linewidth]{HM10.png}
    \caption{Módulo bluetooth HM-10}
    \label{fig:HM10}
\end{figure}

\subsubsection{MAX30102}
En cuanto a la conexión del módulo Bluetooth con la placa de desarrollo, se realizaron cuatro conexiones principales: alimentación (VCC), referencia a tierra (GND) y dos líneas de comunicación correspondientes a la interfaz serie.
Las líneas de datos se conectan mediante una configuración cruzada entre los pines de transmisión y recepción.
En particular, el pin TX (transmit) del módulo se conecta al pin RX (receive) del microcontrolador, mientras que el pin RX del módulo se conecta al pin TX del microcontrolador permitiendo de esta forma el intercambio de datos. 

\begin{figure}[H]
    \centering
    \includegraphics[width=0.5\linewidth]{MAX30102.png}
    \caption{Sensor MAX30102}
    \label{fig: MAX30102}
\end{figure}
